\documentclass[11pt,a4paper]{article}
% preamble.tex — estilo común de todos los apuntes Froth.
% Cada apunte hace % preamble.tex — estilo común de todos los apuntes Froth.
% Cada apunte hace % preamble.tex — estilo común de todos los apuntes Froth.
% Cada apunte hace \input{preamble} justo después de \documentclass.
\usepackage[margin=2.4cm]{geometry}
\usepackage{xcolor}
\definecolor{litblue}{RGB}{31,79,121}
\definecolor{litgreen}{RGB}{27,94,32}
\definecolor{litorange}{RGB}{230,81,0}
\usepackage{parskip}
\usepackage{titlesec}
\titleformat{\section}{\Large\bfseries\color{litblue}}{\thesection}{0.6em}{}
\titleformat{\subsection}{\large\bfseries\color{litblue}}{\thesubsection}{0.6em}{}
\usepackage{tcolorbox}
\usepackage{fancyhdr}
\pagestyle{fancy}
\fancyhf{}
\fancyhead[L]{\small\color{litblue}Froth — Apuntes de ML}
\fancyhead[R]{\small\thepage}
\renewcommand{\headrulewidth}{0.4pt}
\usepackage[colorlinks=true,linkcolor=litblue,urlcolor=litblue]{hyperref}

% Cabecera de cada apunte: \cabecera{numero}{titulo}
\newcommand{\cabecera}[2]{%
  \begin{tcolorbox}[colback=litblue,coltext=white,colframe=litblue,arc=2mm]
    {\Large\bfseries Apunte #1 — #2}\\[3pt]
    {\small Proyecto Froth · aprender Machine Learning construyendo}
  \end{tcolorbox}\vspace{0.8em}}

% Cajas temáticas
\newtcolorbox{concepto}[1]{colback=blue!4,colframe=litblue,title={Concepto: #1},arc=1.5mm}
\newtcolorbox{practica}{colback=green!5,colframe=litgreen,title={Pruébalo tú (teclado en mano)},arc=1.5mm}
\newtcolorbox{ojo}{colback=orange!8,colframe=litorange,title={Ojo / error típico},arc=1.5mm}
\newtcolorbox{resumen}{colback=gray!8,colframe=gray!60,title={En una frase},arc=1.5mm}

\newcommand{\codigo}[1]{\texttt{#1}}
 justo después de \documentclass.
\usepackage[margin=2.4cm]{geometry}
\usepackage{xcolor}
\definecolor{litblue}{RGB}{31,79,121}
\definecolor{litgreen}{RGB}{27,94,32}
\definecolor{litorange}{RGB}{230,81,0}
\usepackage{parskip}
\usepackage{titlesec}
\titleformat{\section}{\Large\bfseries\color{litblue}}{\thesection}{0.6em}{}
\titleformat{\subsection}{\large\bfseries\color{litblue}}{\thesubsection}{0.6em}{}
\usepackage{tcolorbox}
\usepackage{fancyhdr}
\pagestyle{fancy}
\fancyhf{}
\fancyhead[L]{\small\color{litblue}Froth — Apuntes de ML}
\fancyhead[R]{\small\thepage}
\renewcommand{\headrulewidth}{0.4pt}
\usepackage[colorlinks=true,linkcolor=litblue,urlcolor=litblue]{hyperref}

% Cabecera de cada apunte: \cabecera{numero}{titulo}
\newcommand{\cabecera}[2]{%
  \begin{tcolorbox}[colback=litblue,coltext=white,colframe=litblue,arc=2mm]
    {\Large\bfseries Apunte #1 — #2}\\[3pt]
    {\small Proyecto Froth · aprender Machine Learning construyendo}
  \end{tcolorbox}\vspace{0.8em}}

% Cajas temáticas
\newtcolorbox{concepto}[1]{colback=blue!4,colframe=litblue,title={Concepto: #1},arc=1.5mm}
\newtcolorbox{practica}{colback=green!5,colframe=litgreen,title={Pruébalo tú (teclado en mano)},arc=1.5mm}
\newtcolorbox{ojo}{colback=orange!8,colframe=litorange,title={Ojo / error típico},arc=1.5mm}
\newtcolorbox{resumen}{colback=gray!8,colframe=gray!60,title={En una frase},arc=1.5mm}

\newcommand{\codigo}[1]{\texttt{#1}}
 justo después de \documentclass.
\usepackage[margin=2.4cm]{geometry}
\usepackage{xcolor}
\definecolor{litblue}{RGB}{31,79,121}
\definecolor{litgreen}{RGB}{27,94,32}
\definecolor{litorange}{RGB}{230,81,0}
\usepackage{parskip}
\usepackage{titlesec}
\titleformat{\section}{\Large\bfseries\color{litblue}}{\thesection}{0.6em}{}
\titleformat{\subsection}{\large\bfseries\color{litblue}}{\thesubsection}{0.6em}{}
\usepackage{tcolorbox}
\usepackage{fancyhdr}
\pagestyle{fancy}
\fancyhf{}
\fancyhead[L]{\small\color{litblue}Froth — Apuntes de ML}
\fancyhead[R]{\small\thepage}
\renewcommand{\headrulewidth}{0.4pt}
\usepackage[colorlinks=true,linkcolor=litblue,urlcolor=litblue]{hyperref}

% Cabecera de cada apunte: \cabecera{numero}{titulo}
\newcommand{\cabecera}[2]{%
  \begin{tcolorbox}[colback=litblue,coltext=white,colframe=litblue,arc=2mm]
    {\Large\bfseries Apunte #1 — #2}\\[3pt]
    {\small Proyecto Froth · aprender Machine Learning construyendo}
  \end{tcolorbox}\vspace{0.8em}}

% Cajas temáticas
\newtcolorbox{concepto}[1]{colback=blue!4,colframe=litblue,title={Concepto: #1},arc=1.5mm}
\newtcolorbox{practica}{colback=green!5,colframe=litgreen,title={Pruébalo tú (teclado en mano)},arc=1.5mm}
\newtcolorbox{ojo}{colback=orange!8,colframe=litorange,title={Ojo / error típico},arc=1.5mm}
\newtcolorbox{resumen}{colback=gray!8,colframe=gray!60,title={En una frase},arc=1.5mm}

\newcommand{\codigo}[1]{\texttt{#1}}

\begin{document}
\cabecera{36}{TextRank: PageRank sobre frases (el hub de las frases)}

\section{La idea: reciclar el algoritmo de Google}

PageRank ordenaba la web con una idea circular pero resoluble: \emph{una página importa
si páginas importantes la enlazan}. TextRank (Mihalcea \& Tarau, 2004) cambia páginas por
frases y enlaces por similitud: \emph{una frase importa si frases importantes se le
parecen}. Es EXACTAMENTE tu concepto de hub de la Fase 4, aplicado dentro de cada cluster.

\begin{concepto}{el grafo de frases y el paseo aleatorio}
Nodos = las frases del cluster; peso de la arista = coseno entre sus vectores SPECTER.
PageRank se calcula con \textbf{iteración de potencias}: imagina un lector saltando de
frase en frase, prefiriendo saltos hacia frases similares; con probabilidad $1-d$
($d=0.85$) se teletransporta a cualquiera. Tras iterar
\[
r \leftarrow \frac{1-d}{N} + d \, W r
\]
unas 50 veces, $r$ converge: el tiempo que el lector pasa en cada frase ES su importancia.
En \texttt{froth/review.py} son 4 líneas de numpy --- sin librerías nuevas.
\end{concepto}

\section{Centroide vs TextRank: no miden lo mismo}

\begin{itemize}
  \item \textbf{Centroide} premia lo \emph{típico}: la frase más parecida al promedio.
  \item \textbf{TextRank} premia lo \emph{respaldado}: la frase con más ``votos'' de
        consenso, aunque el consenso venga de un subgrupo denso.
\end{itemize}

En tu corpus (¡corre \texttt{python -m froth.review} y júzgalo!):

\begin{itemize}
  \item Cluster USGS: AMBOS eligen ``Minerals Yearbook --- These annual publications
        review...'' --- cuando dos métricas independientes coinciden, esa frase ES el tema.
  \item Cluster lepidolita: TextRank prefirió \emph{``...the lepidolite flotation process
        and its challenges are summarized, including poor selectivity...''} --- frases de
        REVIEW (resumen los retos del campo) reciben muchos votos porque tocan a todas.
        Para un literature review, eso es una virtud.
  \item Cluster geología: TextRank se fue a depósitos exógenos --- un subgrupo denso de
        frases mineralógicas se vota entre sí y captura el ranking. El centroide, anclado
        a los vectores de los PAPERS, resiste mejor esa trampa.
\end{itemize}

\begin{ojo}
Los scores no son comparables entre métodos: el coseno vive en 0--1; PageRank reparte una
probabilidad total de 1 entre N frases (por eso ves 0.008). Compara RANKINGS, no números.
Y ninguno de los dos evita la redundancia --- eso es MMR (7.3).
\end{ojo}

\begin{resumen}
TextRank = PageRank sobre el grafo de similitud de frases, resuelto por iteración de
potencias. Premia el consenso (bueno para frases-resumen), pero un subgrupo denso puede
capturarlo; el centroide premia lo típico. El experto decide cuál lee mejor --- o se
combinan.
\end{resumen}

\end{document}
